\documentclass[11pt]{amsart}

\usepackage[T1]{fontenc}
\usepackage{lmodern}
\usepackage{microtype}
\usepackage{mathtools,amssymb,amsthm}
\usepackage{verbatim}
\usepackage[hidelinks]{hyperref}

\hypersetup{
  pdftitle={A 24-Permutation Hamiltonian Cycle of P(4) That Does Not Extend
  to a u-Cycle}
}

\newtheorem{theorem}{Theorem}
\newtheorem{lemma}[theorem]{Lemma}

\DeclareMathOperator{\stt}{st}
\DeclareMathOperator{\Ham}{Ham}
\DeclareMathOperator{\Eul}{Eul}

\title[A Hamiltonian cycle of $P(4)$ that does not extend to a u-cycle]
{A 24-Permutation Hamiltonian Cycle of $P(4)$ That Does Not Extend to a
u-Cycle}

\date{}

\subjclass[2020]{05A05, 05C45, 05C20, 68R15}
\keywords{universal cycle, u-cycle for permutations, graph of overlapping
permutations, Hamiltonian cycle, linear extension, counterexample}

\begin{document}

\begin{abstract}
Let $P(n)$ be the digraph of overlapping permutations: its vertices are the
$n!$ permutations of $\{1,\dots,n\}$, with an arc $x_1\cdots x_n\to
y_1\cdots y_n$ when $x_2\cdots x_n$ and $y_1\cdots y_{n-1}$ are
order-isomorphic. Chung, Diaconis and Graham introduced a procedure turning a
Hamiltonian cycle of $P(n)$ into a universal cycle for $n$-permutations,
provided the strict inequalities the cycle forces on the letters have no
directed cycle. Kitaev and Qiu report that the authors of that paper believe
this to be the case for \emph{any} Hamiltonian cycle of $P(n)$, and in 2026
restate the belief as a conjecture. We exhibit a directed Hamiltonian cycle of
$P(4)$ whose forced inequalities contain the directed cycle
$u_2<u_5<u_7<u_9<u_{12}<u_{14}<u_{17}<u_{20}<u_{23}<u_2$. It therefore admits
no linear extension and does not extend to a u-cycle, so the universally
quantified statement as those two papers report it fails already at $n=4$. The
proof is a $24$-entry table, $24$ arc checks and nine readings of a window, and
needs no computation. We did not obtain the 1992 paper itself and make no claim
about what it states: what is refuted is the statement as reported in 2024 and
conjectured in 2026. The existence of u-cycles for $n$-permutations, and
Hurlbert's construction to which Chung, Diaconis and Graham refer, are
untouched.
\end{abstract}

\maketitle

\section{The statement that is refuted}

For a word $w$ of distinct letters, $\stt(w)$ denotes its standardisation: the
permutation of $\{1,\dots,|w|\}$ order-isomorphic to $w$. The digraph is
defined by Kitaev and Qiu \cite{KQ2024} as follows; we quote the \LaTeX{}
source of their e-print \texttt{arXiv:2408.05984v1} verbatim, inserting only
line breaks for width.

{\footnotesize\begin{verbatim}
the vertex set of $P(n)$ is the set of all $n!$ permutations of
$\{1,2,\ldots,n\}$, and there is an edge $x_1x_2\cdots x_n\rightarrow
y_1y_2\cdots y_n$ if and only if, for $2\leq i<j\leq n$, $x_i<x_j$ if and
only if $y_{i-1}<y_{j-1}$.
\end{verbatim}}

\noindent
Equivalently, $p\to q$ is an arc of $P(n)$ if and only if
$\stt(p_2\cdots p_n)=\stt(q_1\cdots q_{n-1})$.

Given a directed Hamiltonian cycle $C=(p_0,\dots,p_{N-1})$ of $P(n)$, where
$N=n!$, the Chung--Diaconis--Graham procedure \cite{CDG1992} introduces $N$
undetermined letters $u_0,\dots,u_{N-1}$ placed on a circle and requires, for
every $i$, that the \emph{window}
\[
 (u_i,u_{i+1},\dots,u_{i+n-1}),
 \qquad\text{indices mod }N,
\]
be order-isomorphic to $p_i$. Any two positions lying in a common window
therefore have their relative order forced. Write $D(C)$ for the digraph on
$\{0,\dots,N-1\}$ whose arcs $a\to b$ are the forced strict inequalities
$u_a<u_b$. If $D(C)$ is acyclic then any linear extension of it gives a u-cycle
for $n$-permutations; if $D(C)$ has a directed cycle then no assignment of
letters realises all $N$ windows and $C$ yields nothing.

The belief we refute is reported later in the same file. Quoted exactly,
including the source's own grammatical slip ``believe that this to be the
case'':

{\footnotesize\begin{verbatim}
Even though obtaining a Hamiltonian cycle $C$ in $P(n)$ is straightforward,
arguing that the implied ordering on the values is in fact a partial order,
i.e.\ has no cycles, is a difficult task. While the authors of \cite{CDG1992}
believe that this to be the case for {\em any} Hamiltonian cycle in $P(n)$,
they refer to \cite{Hurlbert} where a particular construction of $C$ is
offered that guarantees the ordering be a partial order.
\end{verbatim}}

\noindent
The load-bearing word is the quantifier \texttt{\{\char`\\em any\}}. The same
two authors restate it as an open conjecture in 2026 \cite{KQ2026}
(\texttt{arXiv:2603.01005v1}); we quote that sentence \emph{together with its
antecedent}, because the antecedent is what fixes the class quantified over.

{\footnotesize\begin{verbatim}
Any Eulerian cycle in a graph formed by clusters can be extended to a
Hamiltonian cycle in the graph of overlapping permutations (since each edge
corresponds to exactly one permutation and we know this permutation). At
least some of these Hamiltonian cycles (possibly all, which is conjectured),
can be extended to u-cycles for permutations via linear extensions of
partially ordered sets as described in~\cite{CDG1992}.
\end{verbatim}}

\begin{lemma}\label{lem:bij}
Let $H_n$ be the multidigraph with vertex set $S_{n-1}$ in which each
$p\in S_n$ is the edge $\stt(p_1\cdots p_{n-1})\to\stt(p_2\cdots p_n)$. Then
$P(n)$ is the line digraph of $H_n$, and the cyclic orderings of $S_n$ whose
consecutive entries are arcs of $P(n)$ are exactly the Eulerian circuits of
$H_n$. In particular $\Ham(P(n))\leftrightarrow\Eul(H_n)$ is a bijection, not
an inclusion.
\end{lemma}

\begin{proof}
By definition $p\to q$ is an arc of $P(n)$ if and only if
$\stt(p_2\cdots p_n)=\stt(q_1\cdots q_{n-1})$, that is, if and only if the head
of the edge $p$ of $H_n$ equals the tail of the edge $q$. A cyclic ordering of
all of $S_n$ with consecutive entries adjacent in $P(n)$ is therefore a closed
walk in $H_n$ using every edge exactly once, and conversely.
\end{proof}

Hence the 2026 phrase ``possibly all'' and the 2024 quantifier ``\emph{any}
Hamiltonian cycle in $P(n)$'' quantify over the same set, and the Eulerian
route of \cite{KQ2026} is a reparametrisation of the whole search space rather
than a construction landing in a distinguished subclass. What is refuted below
is the statement as both sources give it.

\medskip\noindent\textbf{Provenance of the quotations, and what they do not
fix.} The paper \cite{CDG1992} is not open access and we did not read it; the
belief is quoted only through the two citing papers above, whose \texttt{v1}
e-print \LaTeX{} sources we read directly, and the quoted spans are exact
substrings of those files. We therefore do not claim that Chung, Diaconis and
Graham ever asserted the universally quantified statement, still less that they
posed it as a mathematical assertion rather than an expectation: \cite{KQ2024}
reports a belief at second hand, and its own observation that \cite{CDG1992}
refer to \cite{Hurlbert} for a construction carrying a guarantee is consistent
with less having been claimed there. The one firm anchor is the 2026
parenthetical ``possibly all, which is conjectured''. A reader who holds that no
universal assertion was ever posed in \cite{CDG1992} should read what follows as
settling that conjecture of \cite{KQ2026} in the negative and as saying nothing
about \cite{CDG1992}. Neither published journal version of \cite{KQ2024,KQ2026}
was obtainable, so we cannot exclude a copy-edit having altered the quoted
sentences; see Section~\ref{sec:open}.

\medskip\noindent\textbf{Adjacent work.} Isaak \cite{Isaak2006} proves that the
Hurlbert--Jackson digraphs $P(n,k)$, whose vertices are the $k$-permutations of
an $n$-set and whose arcs are the $(k+1)$-permutations, are Hamiltonian for
$k\le n-3$, using restricted Eulerian cycles and a line-digraph relation of the
kind used in Lemma~\ref{lem:bij}; that is an existence result on a different
digraph over an $n$-letter alphabet, and it says nothing about acyclicity of the
order forced by a cycle of the order-isomorphism digraph $P(n)$. Asplund and Fox
\cite{AsplundFox} work with the digraph $P(n)$ itself, count its short cycles
and raise the existence and counting question for Hamiltonian cycles in it;
neither their counts nor that question concerns the order forced on the letters,
so they settle nothing about u-cycles. The conjecture quoted above is restated,
not resolved, in \cite{KQ2026}. We claim no priority for what follows.

\section{The counterexample}

\begin{theorem}\label{thm:main}
Let $n=4$, $N=24$, and let $C=(p_0,\dots,p_{23})$ be
\begin{center}\footnotesize
\begin{tabular}{r|cccccccccccc}
$i$   & 0&1&2&3&4&5&6&7&8&9&10&11\\\hline
$p_i$ & 1234&1243&1432&4312&4132&1324&3142&1423&4123&1342&3421&4213
\end{tabular}

\medskip
\begin{tabular}{r|cccccccccccc}
$i$   & 12&13&14&15&16&17&18&19&20&21&22&23\\\hline
$p_i$ & 2134&2341&2413&4231&2314&2143&2431&4321&3214&3241&3412&3124
\end{tabular}
\end{center}
closed by the arc $p_{23}\to p_0$, that is $3124\to1234$. Then $C$ is a
directed Hamiltonian cycle of $P(4)$ and $D(C)$ contains the directed cycle
\begin{equation}\label{eq:chain}
 u_2<u_5<u_7<u_9<u_{12}<u_{14}<u_{17}<u_{20}<u_{23}<u_2 .
\end{equation}
Hence the implied ordering is not a partial order, $C$ has no linear extension,
and $C$ does not extend to a u-cycle for $4$-permutations. Consequently the
universally quantified statement reported in \cite{KQ2024} and conjectured in
\cite{KQ2026} fails already at $n=4$.
\end{theorem}

\begin{proof}
\emph{Every arc holds.} With the rule
$p\to q\iff\stt(p_2p_3p_4)=\stt(q_1q_2q_3)$:
\begin{center}\small
\begin{tabular}{@{}l@{\ }l@{\qquad}l@{\ }l@{}}
$0\to1$ & $\stt(234)=123=\stt(124)$ & $12\to13$ & $\stt(134)=123=\stt(234)$\\
$1\to2$ & $\stt(243)=132=\stt(143)$ & $13\to14$ & $\stt(341)=231=\stt(241)$\\
$2\to3$ & $\stt(432)=321=\stt(431)$ & $14\to15$ & $\stt(413)=312=\stt(423)$\\
$3\to4$ & $\stt(312)=312=\stt(413)$ & $15\to16$ & $\stt(231)=231=\stt(231)$\\
$4\to5$ & $\stt(132)=132=\stt(132)$ & $16\to17$ & $\stt(314)=213=\stt(214)$\\
$5\to6$ & $\stt(324)=213=\stt(314)$ & $17\to18$ & $\stt(143)=132=\stt(243)$\\
$6\to7$ & $\stt(142)=132=\stt(142)$ & $18\to19$ & $\stt(431)=321=\stt(432)$\\
$7\to8$ & $\stt(423)=312=\stt(412)$ & $19\to20$ & $\stt(321)=321=\stt(321)$\\
$8\to9$ & $\stt(123)=123=\stt(134)$ & $20\to21$ & $\stt(214)=213=\stt(324)$\\
$9\to10$ & $\stt(342)=231=\stt(342)$ & $21\to22$ & $\stt(241)=231=\stt(341)$\\
$10\to11$ & $\stt(421)=321=\stt(421)$ & $22\to23$ & $\stt(412)=312=\stt(312)$\\
$11\to12$ & $\stt(213)=213=\stt(213)$ & $23\to0$ & $\stt(124)=123=\stt(123)$
\end{tabular}
\end{center}
The last entry is the wrap arc, so this is a cycle and not a path. No $p_i$
equals $p_{i+1}$, so neither of the two self-loops of $P(4)$ --- which are
$1234\to1234$ and $4321\to4321$, and which no Hamiltonian cycle may use ---
lies on $C$.

\emph{It is Hamiltonian.} Group the entries by first letter:
\begin{center}\small
\begin{tabular}{c|l}
first letter & suffixes occurring\\\hline
$1$ & $234,\ 243,\ 432,\ 324,\ 423,\ 342$\\
$2$ & $134,\ 341,\ 413,\ 314,\ 143,\ 431$\\
$3$ & $142,\ 421,\ 214,\ 241,\ 412,\ 124$\\
$4$ & $312,\ 132,\ 123,\ 213,\ 231,\ 321$
\end{tabular}
\end{center}
Each row lists all six permutations of the three remaining letters, and
$4\cdot6=24=|S_4|$, so the $24$ entries are distinct and exhaust $S_4$.

\emph{The contradiction.} Window $i$ occupies the cyclic positions
$i,i+1,i+2,i+3$ and is order-isomorphic to $p_i$, so the rank of $u_{i+k-1}$
inside window $i$ is the $k$-th letter of $p_i$. Nine windows suffice, each
contributing the inequality between the two letters of consecutive rank shown:
\begin{center}\small
\begin{tabular}{c|c|c|l|c}
window $i$ & $p_i$ & positions & ranks & forced\\\hline
$2$ & $1432$ & $(2,3,4,5)$ & $u_2{=}1,\ u_3{=}4,\ u_4{=}3,\ u_5{=}2$ & $u_2<u_5$\\
$5$ & $1324$ & $(5,6,7,8)$ & $u_5{=}1,\ u_6{=}3,\ u_7{=}2,\ u_8{=}4$ & $u_5<u_7$\\
$6$ & $3142$ & $(6,7,8,9)$ & $u_6{=}3,\ u_7{=}1,\ u_8{=}4,\ u_9{=}2$ & $u_7<u_9$\\
$9$ & $1342$ & $(9,10,11,12)$ & $u_9{=}1,\ u_{10}{=}3,\ u_{11}{=}4,\ u_{12}{=}2$ & $u_9<u_{12}$\\
$12$ & $2134$ & $(12,13,14,15)$ & $u_{12}{=}2,\ u_{13}{=}1,\ u_{14}{=}3,\ u_{15}{=}4$ & $u_{12}<u_{14}$\\
$14$ & $2413$ & $(14,15,16,17)$ & $u_{14}{=}2,\ u_{15}{=}4,\ u_{16}{=}1,\ u_{17}{=}3$ & $u_{14}<u_{17}$\\
$17$ & $2143$ & $(17,18,19,20)$ & $u_{17}{=}2,\ u_{18}{=}1,\ u_{19}{=}4,\ u_{20}{=}3$ & $u_{17}<u_{20}$\\
$20$ & $3214$ & $(20,21,22,23)$ & $u_{20}{=}3,\ u_{21}{=}2,\ u_{22}{=}1,\ u_{23}{=}4$ & $u_{20}<u_{23}$\\
$23$ & $3124$ & $(23,0,1,2)$ & $u_{23}{=}3,\ u_0{=}1,\ u_1{=}2,\ u_2{=}4$ & $u_{23}<u_2$
\end{tabular}
\end{center}
The last row is a wrap-around window. Concatenating the nine forced
inequalities gives \eqref{eq:chain}, a cycle of \emph{strict} inequalities, so
no values whatsoever can be assigned to $u_0,\dots,u_{23}$.
\end{proof}

\section{The reading of the source}

\noindent\textbf{The wrap-around windows, pinned to published
output.} Everything above rests on reading all $N$ windows, including the $n-1$
that wrap around; omitting those is the single error that would manufacture a
false ``acyclic''. That the wrap-around reading is the source's own is settled
by the source's own worked example at $n=3$ in
\texttt{arXiv:2408.05984v1}: for the Hamiltonian cycle
\[
 132\to312\to123\to231\to321\to213\to132
\]
of $P(3)$ they publish the u-cycle word $U_3=142342$. Read as a
\emph{circular} word of length $6$, its six windows of length $3$ are
\[
 (1,4,2),\ (4,2,3),\ (2,3,4),\ (3,4,2),\ (4,2,1),\ (2,1,4)
 \;\longrightarrow\;
 132,\ 312,\ 123,\ 231,\ 321,\ 213,
\]
their cycle in order, \emph{including} the two wrap-around windows $(4,2,1)$
and $(2,1,4)$. Applying the forced-order construction to their cycle returns
the transitive reduction
$\{a\to c,\ a\to f,\ c\to d,\ f\to d,\ d\to b,\ d\to e\}$ on the six positions
$a,\dots,f$, which is their published Hasse diagram, and $142342$ is one of its
linear extensions. By contrast, dropping the wrap-around windows makes the
forced order of the cycle of Theorem~\ref{thm:main} acyclic --- which is exactly
the failure mode this control excludes.

\medskip\noindent\textbf{Independence of the reading.} Theorem~\ref{thm:main}
does not depend on which of the two natural readings of ``the implied
ordering'' is taken. The source's literal all-ordered-pairs reading gives $72$
forced strict inequalities on the exhibited cycle, with no antisymmetry
violation; the consecutive-rank reading used in the proof gives $46$. Both
contain all nine links of \eqref{eq:chain}.

\section{What is not settled}\label{sec:open}

Only $n=4$ is settled. Whether every $n\ge5$ admits a Hamiltonian cycle of
$P(n)$ whose forced order is cyclic is not settled here, and we report no
computation at $n\ge5$. Neither do we report how many Hamiltonian cycles of
$P(4)$ have cyclic forced order: no such count is claimed anywhere in this
paper. The arc counts $72$ and $46$ are about the exhibited cycle alone.

\medskip\noindent\textbf{Sources not read.} The paper \cite{CDG1992} itself (no
open-access copy
exists; publisher access refused), Hurlbert's 1990 Rutgers thesis
\cite{Hurlbert} (not online, no repository record; it is the source of the
existential construction, which this paper does not contradict), and the survey
\cite{JSH2009} (full text refused at the only location we found). Neither
paywalled journal version of \cite{KQ2024,KQ2026} was obtainable, so the
published-versus-preprint text was never diffed; the risk is low, since
\cite{KQ2026} independently restates the conjecture as open in 2026.

\medskip\noindent\textbf{What is not claimed.} U-cycles for $n$-permutations
exist for every $n$, and every published construction that actually produces
one still works; Hurlbert's construction \cite{Hurlbert}, to which
\cite{CDG1992} refer, is untouched. What fails is only the
shortcut of starting from an \emph{arbitrary} Hamiltonian cycle. By
Lemma~\ref{lem:bij}, ``take an Eulerian cycle of the cluster graph and lift
it'' \cite{KQ2026} is not an alternative guarantee: it generates every
Hamiltonian cycle of $P(n)$, the one exhibited above included, and
\cite{KQ2026} itself claims no more for it than that ``at least some'' of the
resulting cycles work.

\section{Reproduction and verification}

Everything in Theorem~\ref{thm:main} is checkable from the printed tables with
no computer. The program \texttt{verify.py} accompanying this paper reads the
$24$-entry cycle printed above, together with the published $n=3$ objects of
\cite{KQ2024}, and re-derives: that the cycle is a Hamiltonian cycle of $P(4)$
using neither self-loop; that each of the nine links of \eqref{eq:chain} is
forced by the window named in the table; that the forced digraph is cyclic under
both readings of the source; the arc counts $46$ and $72$ and the absence of
antisymmetry violations; and that Kitaev and Qiu's own cycle, word
$U_3=142342$ and Hasse diagram are reproduced by the convention used here, while
dropping the wrap-around windows would make the exhibited cycle look acyclic. It
reports $73$ checks, all passing, in exact integer arithmetic, using only the
Python standard library.

The recorded output covers more than this paper claims, and the surplus is not
claimed here: a second $24$-permutation Hamiltonian cycle of $P(4)$ with cyclic
forced order, which this paper does not print; exhaustive verdicts at $n=3$ and
$n=2$; and a total number of Hamiltonian cycles of $P(4)$ obtained from an
integer Matrix--Tree cofactor with published de Bruijn control counts. The
program re-derives nothing about $n\ge5$. The verbatim quotations above are
facts about external e-print files and are not checked by the program.

\begin{thebibliography}{9}

\bibitem{AsplundFox}
J.~Asplund and N.~Fox,
\emph{Enumerating cycles in the graph of overlapping permutations},
Discrete Math. \textbf{341} (2018), no.~2, 427--438.
\href{https://arxiv.org/abs/1609.02210}{arXiv:1609.02210}.

\bibitem{CDG1992}
F.~Chung, P.~Diaconis, and R.~Graham,
\emph{Universal cycles for combinatorial structures},
Discrete Math. \textbf{110} (1992), 43--59.
\href{https://doi.org/10.1016/0012-365X(92)90699-G}
{doi:10.1016/0012-365X(92)90699-G}.
Not open access; we did not read it, and quote its belief only through
\cite{KQ2024,KQ2026}.

\bibitem{Hurlbert}
G.~Hurlbert,
Ph.D. thesis, Rutgers University, 1990.
Cited as \texttt{\char`\\cite\{Hurlbert\}} by \cite{KQ2024} for the
construction that does carry a guarantee; we could not obtain it, and record
only the bibliographic data \cite{KQ2024} gives.

\bibitem{Isaak2006}
G.~Isaak,
\emph{Hamiltonicity of digraphs for universal cycles of permutations},
European J. Combin. \textbf{27} (2006), no.~6, 801--805.
\href{https://doi.org/10.1016/j.ejc.2005.05.007}
{doi:10.1016/j.ejc.2005.05.007}.

\bibitem{JSH2009}
B.~Jackson, B.~Stevens, and G.~Hurlbert,
\emph{Research problems on Gray codes and universal cycles},
Discrete Math. \textbf{309} (2009), 5341--5348.
Full text not obtainable to us.

\bibitem{KQ2024}
S.~Kitaev and D.~Qiu,
\emph{On a family of universal cycles for multi-dimensional permutations},
\href{https://arxiv.org/abs/2408.05984}{arXiv:2408.05984v1}, 2024.
The passages quoted above are exact substrings of the \LaTeX{} source of that
e-print.

\bibitem{KQ2026}
S.~Kitaev and D.~Qiu,
\emph{On shortening universal words for multi-dimensional permutations},
\href{https://arxiv.org/abs/2603.01005}{arXiv:2603.01005v1}, 2026.
The passage quoted above is an exact substring of the \LaTeX{} source of that
e-print.

\end{thebibliography}

\end{document}
